% Options for packages loaded elsewhere
\PassOptionsToPackage{unicode}{hyperref}
\PassOptionsToPackage{hyphens}{url}
\PassOptionsToPackage{dvipsnames,svgnames,x11names}{xcolor}
\documentclass[
  11pt,
  a4paper,
]{article}
\usepackage{xcolor}
\usepackage[margin=1in]{geometry}
\usepackage{amsmath,amssymb}
\setcounter{secnumdepth}{-\maxdimen} % remove section numbering
\usepackage{iftex}
\ifPDFTeX
  \usepackage[T1]{fontenc}
  \usepackage[utf8]{inputenc}
  \usepackage{textcomp} % provide euro and other symbols
\else % if luatex or xetex
  \usepackage{unicode-math} % this also loads fontspec
  \defaultfontfeatures{Scale=MatchLowercase}
  \defaultfontfeatures[\rmfamily]{Ligatures=TeX,Scale=1}
\fi
\usepackage{lmodern}
\ifPDFTeX\else
  % xetex/luatex font selection
  \ifLuaTeX
    \usepackage{luaotfload}
    \directlua{luaotfload.add_fallback("mainfontfallback",{
      "NotoSerifBengali:mode=harf"
    })}
  \fi
  \setmainfont[,RawFeature={fallback=mainfontfallback}]{Latin Modern Roman}
  \ifLuaTeX
    \usepackage{luaotfload}
    \directlua{luaotfload.add_fallback("monofontfallback",{
      "NotoSansBengali:mode=harf"
    })}
  \fi
  \setmonofont[,RawFeature={fallback=monofontfallback}]{DejaVu Sans Mono}
\fi
% Use upquote if available, for straight quotes in verbatim environments
\IfFileExists{upquote.sty}{\usepackage{upquote}}{}
\IfFileExists{microtype.sty}{% use microtype if available
  \usepackage[]{microtype}
  \UseMicrotypeSet[protrusion]{basicmath} % disable protrusion for tt fonts
}{}
\makeatletter
\@ifundefined{KOMAClassName}{% if non-KOMA class
  \IfFileExists{parskip.sty}{%
    \usepackage{parskip}
  }{% else
    \setlength{\parindent}{0pt}
    \setlength{\parskip}{6pt plus 2pt minus 1pt}}
}{% if KOMA class
  \KOMAoptions{parskip=half}}
\makeatother
\usepackage{color}
\usepackage{fancyvrb}
\newcommand{\VerbBar}{|}
\newcommand{\VERB}{\Verb[commandchars=\\\{\}]}
\DefineVerbatimEnvironment{Highlighting}{Verbatim}{commandchars=\\\{\}}
% Add ',fontsize=\small' for more characters per line
\usepackage{framed}
\definecolor{shadecolor}{RGB}{248,248,248}
\newenvironment{Shaded}{\begin{snugshade}}{\end{snugshade}}
\newcommand{\AlertTok}[1]{\textcolor[rgb]{0.94,0.16,0.16}{#1}}
\newcommand{\AnnotationTok}[1]{\textcolor[rgb]{0.56,0.35,0.01}{\textbf{\textit{#1}}}}
\newcommand{\AttributeTok}[1]{\textcolor[rgb]{0.13,0.29,0.53}{#1}}
\newcommand{\BaseNTok}[1]{\textcolor[rgb]{0.00,0.00,0.81}{#1}}
\newcommand{\BuiltInTok}[1]{#1}
\newcommand{\CharTok}[1]{\textcolor[rgb]{0.31,0.60,0.02}{#1}}
\newcommand{\CommentTok}[1]{\textcolor[rgb]{0.56,0.35,0.01}{\textit{#1}}}
\newcommand{\CommentVarTok}[1]{\textcolor[rgb]{0.56,0.35,0.01}{\textbf{\textit{#1}}}}
\newcommand{\ConstantTok}[1]{\textcolor[rgb]{0.56,0.35,0.01}{#1}}
\newcommand{\ControlFlowTok}[1]{\textcolor[rgb]{0.13,0.29,0.53}{\textbf{#1}}}
\newcommand{\DataTypeTok}[1]{\textcolor[rgb]{0.13,0.29,0.53}{#1}}
\newcommand{\DecValTok}[1]{\textcolor[rgb]{0.00,0.00,0.81}{#1}}
\newcommand{\DocumentationTok}[1]{\textcolor[rgb]{0.56,0.35,0.01}{\textbf{\textit{#1}}}}
\newcommand{\ErrorTok}[1]{\textcolor[rgb]{0.64,0.00,0.00}{\textbf{#1}}}
\newcommand{\ExtensionTok}[1]{#1}
\newcommand{\FloatTok}[1]{\textcolor[rgb]{0.00,0.00,0.81}{#1}}
\newcommand{\FunctionTok}[1]{\textcolor[rgb]{0.13,0.29,0.53}{\textbf{#1}}}
\newcommand{\ImportTok}[1]{#1}
\newcommand{\InformationTok}[1]{\textcolor[rgb]{0.56,0.35,0.01}{\textbf{\textit{#1}}}}
\newcommand{\KeywordTok}[1]{\textcolor[rgb]{0.13,0.29,0.53}{\textbf{#1}}}
\newcommand{\NormalTok}[1]{#1}
\newcommand{\OperatorTok}[1]{\textcolor[rgb]{0.81,0.36,0.00}{\textbf{#1}}}
\newcommand{\OtherTok}[1]{\textcolor[rgb]{0.56,0.35,0.01}{#1}}
\newcommand{\PreprocessorTok}[1]{\textcolor[rgb]{0.56,0.35,0.01}{\textit{#1}}}
\newcommand{\RegionMarkerTok}[1]{#1}
\newcommand{\SpecialCharTok}[1]{\textcolor[rgb]{0.81,0.36,0.00}{\textbf{#1}}}
\newcommand{\SpecialStringTok}[1]{\textcolor[rgb]{0.31,0.60,0.02}{#1}}
\newcommand{\StringTok}[1]{\textcolor[rgb]{0.31,0.60,0.02}{#1}}
\newcommand{\VariableTok}[1]{\textcolor[rgb]{0.00,0.00,0.00}{#1}}
\newcommand{\VerbatimStringTok}[1]{\textcolor[rgb]{0.31,0.60,0.02}{#1}}
\newcommand{\WarningTok}[1]{\textcolor[rgb]{0.56,0.35,0.01}{\textbf{\textit{#1}}}}
\usepackage{longtable,booktabs,array}
\usepackage{calc} % for calculating minipage widths
% Correct order of tables after \paragraph or \subparagraph
\usepackage{etoolbox}
\makeatletter
\patchcmd\longtable{\par}{\if@noskipsec\mbox{}\fi\par}{}{}
\makeatother
% Allow footnotes in longtable head/foot
\IfFileExists{footnotehyper.sty}{\usepackage{footnotehyper}}{\usepackage{footnote}}
\makesavenoteenv{longtable}
\setlength{\emergencystretch}{3em} % prevent overfull lines
\providecommand{\tightlist}{%
  \setlength{\itemsep}{0pt}\setlength{\parskip}{0pt}}
% paper_preamble.tex — shared LaTeX preamble for the PCMA/SI paper set.
%
% Compiled with lualatex (NOT pdflatex) because several papers contain
% Bengali script: the CLL paper realises one cognitive state into
% English, Bengali and Banglish, so the Bengali output is EVIDENCE and
% cannot be transliterated away without destroying the point being made.
%
% Font selection is NOT done here. pandoc >= 3 has native
% `mainfontfallback` / `monofontfallback` variables that emit the
% luaotfload fallback registration in the correct position relative to
% unicode-math; doing it by hand in this header lands after unicode-math
% and breaks font loading outright. See tools/build_papers.sh.

\usepackage{amsmath}
\usepackage{amssymb}
\usepackage{mathtools}   % \xleftrightarrow etc. used in cll.md §5
\usepackage{booktabs}

% ---------------------------------------------------------------------
% Unicode symbols appearing in PROSE rather than inside $...$.
%
% The papers were written as markdown and read as markdown, so they use
% real characters — "the ω-boost", "α-child", "ΔR coherence" — instead of
% LaTeX macros. Under unicode-math those codepoints are math-mode-only
% and abort the build ("\symbballowed only in math mode"). Mapping each
% to \ensuremath keeps the markdown readable on GitHub AND compilable,
% which is the whole point of keeping one source.
% ---------------------------------------------------------------------
\usepackage{newunicodechar}

% Greek, lower case
\newunicodechar{α}{\ensuremath{\alpha}}
\newunicodechar{β}{\ensuremath{\beta}}
\newunicodechar{γ}{\ensuremath{\gamma}}
\newunicodechar{δ}{\ensuremath{\delta}}
\newunicodechar{ε}{\ensuremath{\varepsilon}}
\newunicodechar{η}{\ensuremath{\eta}}
\newunicodechar{θ}{\ensuremath{\theta}}
\newunicodechar{κ}{\ensuremath{\kappa}}
\newunicodechar{λ}{\ensuremath{\lambda}}
\newunicodechar{μ}{\ensuremath{\mu}}
\newunicodechar{π}{\ensuremath{\pi}}
\newunicodechar{ρ}{\ensuremath{\rho}}
\newunicodechar{σ}{\ensuremath{\sigma}}
\newunicodechar{τ}{\ensuremath{\tau}}
\newunicodechar{φ}{\ensuremath{\varphi}}
\newunicodechar{χ}{\ensuremath{\chi}}
\newunicodechar{ψ}{\ensuremath{\psi}}
\newunicodechar{ω}{\ensuremath{\omega}}
\newunicodechar{ν}{\ensuremath{\nu}}
\newunicodechar{ξ}{\ensuremath{\xi}}
\newunicodechar{ζ}{\ensuremath{\zeta}}
\newunicodechar{ι}{\ensuremath{\iota}}

% Greek, upper case
\newunicodechar{Δ}{\ensuremath{\Delta}}
\newunicodechar{Φ}{\ensuremath{\Phi}}
\newunicodechar{Σ}{\ensuremath{\Sigma}}
\newunicodechar{Ψ}{\ensuremath{\Psi}}
\newunicodechar{Ω}{\ensuremath{\Omega}}
\newunicodechar{Γ}{\ensuremath{\Gamma}}
\newunicodechar{Λ}{\ensuremath{\Lambda}}
\newunicodechar{Θ}{\ensuremath{\Theta}}
\newunicodechar{Π}{\ensuremath{\Pi}}

% Operators and relations
\newunicodechar{×}{\ensuremath{\times}}
\newunicodechar{·}{\ensuremath{\cdot}}
\newunicodechar{±}{\ensuremath{\pm}}
\newunicodechar{−}{\ensuremath{-}}
\newunicodechar{≈}{\ensuremath{\approx}}
\newunicodechar{≤}{\ensuremath{\leq}}
\newunicodechar{≥}{\ensuremath{\geq}}
\newunicodechar{≠}{\ensuremath{\neq}}
\newunicodechar{∈}{\ensuremath{\in}}
\newunicodechar{∂}{\ensuremath{\partial}}
\newunicodechar{∇}{\ensuremath{\nabla}}
\newunicodechar{√}{\ensuremath{\sqrt{}}}
\newunicodechar{∞}{\ensuremath{\infty}}
\newunicodechar{∝}{\ensuremath{\propto}}
\newunicodechar{∑}{\ensuremath{\sum}}
\newunicodechar{∫}{\ensuremath{\int}}
\newunicodechar{ℝ}{\ensuremath{\mathbb{R}}}
\newunicodechar{ℕ}{\ensuremath{\mathbb{N}}}
\newunicodechar{‖}{\ensuremath{\|}}
\newunicodechar{⊕}{\ensuremath{\oplus}}
\newunicodechar{⊗}{\ensuremath{\otimes}}
\newunicodechar{∀}{\ensuremath{\forall}}
\newunicodechar{∃}{\ensuremath{\exists}}
\newunicodechar{¬}{\ensuremath{\neg}}
\newunicodechar{∧}{\ensuremath{\wedge}}
\newunicodechar{∨}{\ensuremath{\vee}}

% Arrows
\newunicodechar{→}{\ensuremath{\rightarrow}}
\newunicodechar{←}{\ensuremath{\leftarrow}}
\newunicodechar{↔}{\ensuremath{\leftrightarrow}}
\newunicodechar{↑}{\ensuremath{\uparrow}}
\newunicodechar{↓}{\ensuremath{\downarrow}}
\newunicodechar{⇒}{\ensuremath{\Rightarrow}}
\newunicodechar{⇐}{\ensuremath{\Leftarrow}}
\newunicodechar{⇔}{\ensuremath{\Leftrightarrow}}

% Super/subscripts written as literal characters
\newunicodechar{²}{\ensuremath{^{2}}}
\newunicodechar{³}{\ensuremath{^{3}}}
\newunicodechar{⁴}{\ensuremath{^{4}}}
\newunicodechar{⁵}{\ensuremath{^{5}}}
\newunicodechar{⁶}{\ensuremath{^{6}}}
\newunicodechar{₀}{\ensuremath{_{0}}}
\newunicodechar{₁}{\ensuremath{_{1}}}
\newunicodechar{₂}{\ensuremath{_{2}}}
\newunicodechar{₃}{\ensuremath{_{3}}}
\newunicodechar{ᵀ}{\ensuremath{^{\mathsf{T}}}}

% Status glyphs. The papers use emoji as status markers in tables
% (implemented / failed / warning). Mapped to typographic equivalents:
% a colour emoji font in lualatex is unreliable, and the MEANING is what
% carries — a reader needs to see "this one failed", not a red cross.
\newunicodechar{✅}{\ensuremath{\checkmark}}
\newunicodechar{✓}{\ensuremath{\checkmark}}
\newunicodechar{❌}{\ensuremath{\times}}
\newunicodechar{✗}{\ensuremath{\times}}
\newunicodechar{⚠}{\textbf{!}}
\newunicodechar{🧩}{\textbf{[piece]}}
\newunicodechar{📋}{\textbf{[list]}}
\newunicodechar{🟡}{\textbf{[partial]}}
\newunicodechar{🔴}{\textbf{[dead]}}
\newunicodechar{🟢}{\textbf{[live]}}
\newunicodechar{▼}{\ensuremath{\blacktriangledown}}
\newunicodechar{►}{\ensuremath{\blacktriangleright}}
\newunicodechar{□}{\ensuremath{\square}}
\newunicodechar{█}{\rule{0.6em}{0.9ex}}
\newunicodechar{•}{\textbullet}
\newunicodechar{†}{\textdagger}
\newunicodechar{‑}{-}
\newunicodechar{‐}{-}
\newunicodechar{…}{\ldots}
\newunicodechar{।}{\textbar}   % danda: rendered by the Bengali fallback in text
\newunicodechar{⁻}{\ensuremath{^{-}}}

% Zero-width formatting codepoints. These are invisible in an editor but
% are real characters to LaTeX: an emoji written as "⚠️" is TWO
% codepoints (U+26A0 U+FE0F), so declaring only the visible one leaves a
% variation selector with no definition and the build dies on a symbol
% nobody can see. Both are mapped to nothing.
\newunicodechar{^^^^fe0f}{}    % VARIATION SELECTOR-16
\newunicodechar{^^^^200d}{}    % ZERO WIDTH JOINER
\newunicodechar{^^^^200b}{}    % ZERO WIDTH SPACE

\usepackage{longtable}
\usepackage{array}

% Comparison matrices and cost ledgers are wide. Let them shrink to the
% text block rather than overflow the margin.
\usepackage{adjustbox}

% Code blocks: wrap long lines instead of running off the page. Several
% papers quote real C++ with long identifiers.
\usepackage{fvextra}
\fvset{breaklines=true,breakanywhere=true,fontsize=\small}

\setlength{\parskip}{0.5em}
\setlength{\parindent}{0pt}

\usepackage{fancyhdr}
\pagestyle{fancy}
\fancyhf{}
\fancyfoot[C]{\thepage}
\fancyhead[L]{\small\itshape Synthetic Intelligence — research lineage}
\fancyhead[R]{\small\itshape irfan@furylogic.com}
\renewcommand{\headrulewidth}{0.2pt}
\setlength{\headheight}{14pt}
\setlength{\headsep}{18pt}
\usepackage{bookmark}
\IfFileExists{xurl.sty}{\usepackage{xurl}}{} % add URL line breaks if available
\urlstyle{same}
\hypersetup{
  pdftitle={The Physics Substrate --- Cavities, Materials, Oscillators and Gating},
  pdfauthor={Irfan Mahir --- Independent Researcher, Dhaka, Bangladesh --- irfan@furylogic.com},
  colorlinks=true,
  linkcolor={black},
  filecolor={Maroon},
  citecolor={Blue},
  urlcolor={Blue},
  pdfcreator={LaTeX via pandoc}}

\title{The Physics Substrate --- Cavities, Materials, Oscillators and
Gating}
\author{Irfan Mahir --- Independent Researcher, Dhaka, Bangladesh ---
irfan@furylogic.com}
\date{}

\begin{document}
\maketitle

{
\hypersetup{linkcolor=}
\setcounter{tocdepth}{2}
\tableofcontents
}
\section{The Physics Substrate --- Cavities, Materials, Oscillators and
Gating}\label{the-physics-substrate-cavities-materials-oscillators-and-gating}

\textbf{Irfan Mahir}\\
\emph{Independent Researcher, Dhaka, Bangladesh}\\
irfan@furylogic.com

\begin{quote}
Layer 1. How a concept becomes a mechanical object and a relation
becomes a material.
\end{quote}

\begin{center}\rule{0.5\linewidth}{0.5pt}\end{center}

\subsection{Contents of This Volume}\label{contents-of-this-volume}

\begin{longtable}[]{@{}
  >{\raggedright\arraybackslash}p{(\linewidth - 4\tabcolsep) * \real{0.3333}}
  >{\raggedright\arraybackslash}p{(\linewidth - 4\tabcolsep) * \real{0.3333}}
  >{\raggedright\arraybackslash}p{(\linewidth - 4\tabcolsep) * \real{0.3333}}@{}}
\toprule\noalign{}
\begin{minipage}[b]{\linewidth}\raggedright
\#
\end{minipage} & \begin{minipage}[b]{\linewidth}\raggedright
Paper
\end{minipage} & \begin{minipage}[b]{\linewidth}\raggedright
File
\end{minipage} \\
\midrule\noalign{}
\endhead
\bottomrule\noalign{}
\endlastfoot
1 & PRM: The Phonic Resonance Matrix --- Concepts as Resonant Cavities
and Inference as Wave Propagation & \texttt{prm.md} \\
2 & CMD: Continuous Morphological Dynamics --- Sensory Grounding and
Thermodynamic Habituation & \texttt{cmd.md} \\
3 & MIMS: Mechanical Impedance and Material Stiffness --- Relation Type
as a Physical Property & \texttt{mims.md} \\
4 & HNCA: Hierarchical Neural Constraint Architecture --- Cognition as
Constraint Satisfaction across Two Coupled Scales & \texttt{hnca.md} \\
5 & Neuro-Topological Dynamics Architecture (NTDA): Translating
Biological Brain Physics into Graph-Rewriting Systems &
\texttt{ntda.md} \\
6 & Topological Salience and Default-Mode Architecture (TSDA): Emulating
Human Intuition and Spontaneous Cognition in Graph-Rewriting Systems &
\texttt{tsda.md} \\
7 & SRG: Synaptic Receptor Gating --- Dual-Port Cavities as an
Alternative to Structural Fission & \texttt{srg.md} \\
\end{longtable}

\begin{center}\rule{0.5\linewidth}{0.5pt}\end{center}

\newpage

\subsection{PRM: The Phonic Resonance Matrix --- Concepts as Resonant
Cavities and Inference as Wave
Propagation}\label{prm-the-phonic-resonance-matrix-concepts-as-resonant-cavities-and-inference-as-wave-propagation}

\begin{quote}
\emph{``Inhibition is not an \texttt{if} statement. It is two waves
arriving out of phase.''}
\end{quote}

\begin{center}\rule{0.5\linewidth}{0.5pt}\end{center}

\subsubsection{Abstract}\label{abstract}

We specify a substrate in which a concept is not an identifier, a token
or an embedding, but a \textbf{physical resonant cavity} with mass,
displacement and velocity. Relations are not pointers but
\textbf{tunnels} that transmit mechanical force between cavities.
Inference is not search but \textbf{wave propagation}, and the answer to
a query is whichever cavity is vibrating hardest once the field has
settled.

The architecture's distinguishing property is that several capabilities
usually requiring explicit algorithms become consequences of the medium.
Inhibition is destructive interference rather than a conditional.
Multi-constraint search is two wavefronts colliding rather than a set
intersection. Structural analogy is sympathetic resonance between
cavities that share a harmonic ratio and no edge at all.

Section 7 reports what was measured and what remains an argument from
mechanism, including one claim in the original formulation that does not
survive scrutiny.

\begin{center}\rule{0.5\linewidth}{0.5pt}\end{center}

\subsubsection{1. The Substrate}\label{the-substrate}

Concepts are \textbf{resonant volumes}. Each cavity carries a mass, a
displacement \(x\), a velocity \(v\), and a natural frequency determined
by its physical parameters. In the acoustic reading, a large cavity has
a low resonant frequency and a small one a high frequency.

Relations are \textbf{physical tunnels} --- channels connecting
cavities, not references. Two polarities exist:

\begin{itemize}
\tightlist
\item
  An \textbf{excitatory} tunnel transmits displacement directly,
  producing sympathetic resonance in the target cavity.
\item
  An \textbf{inhibitory} tunnel inverts phase by \(180^\circ\), so its
  contribution actively cancels vibration at the target.
\end{itemize}

The design commitment is that the state of a concept is its \emph{actual
mechanical state}, not a value standing for one.

\begin{center}\rule{0.5\linewidth}{0.5pt}\end{center}

\subsubsection{2. Dynamics}\label{dynamics}

The tick integrates Hooke's law and Newton's second law across every
tunnel:

\[
F \;=\; k\,(x_a \cdot p - x_b)\;-\;c\,(v_a - v_b)
\]

\[
v \;\mathrel{+}=\; \frac{F}{m}\,\Delta t,
\qquad
x \;\mathrel{+}=\; v\,\Delta t
\]

where \(k\) is tunnel stiffness, \(c\) damping, \(p \in \{+1,-1\}\)
polarity, and forces are applied equal and opposite to the two
endpoints.

A query is an injected driving force rather than a string. Reading the
result means observing which cavities have the largest displacement
after the field settles --- the physical analogue of placing pressure
sensors in the pipes and seeing which are shaking.

\begin{center}\rule{0.5\linewidth}{0.5pt}\end{center}

\subsubsection{3. Why Inhibition Requires No
Logic}\label{why-inhibition-requires-no-logic}

In a conventional implementation, negation is a branch: \emph{if the
relation is inhibitory, subtract instead of add.} Here the polarity term
\(p\) sits inside the force expression, so an inhibitory tunnel produces
a \(180^\circ\)-shifted contribution by arithmetic alone.

The consequence is architectural rather than aesthetic. Because
inhibition is multiplication by \(-1\) rather than a conditional, the
update path contains no branch on relation semantics --- which is what
permits the whole tick to be expressed as scalar arithmetic and
saturating nonlinearities, with no control flow that depends on what a
relation \emph{means}.

\begin{center}\rule{0.5\linewidth}{0.5pt}\end{center}

\subsubsection{4. Emergent Capabilities}\label{emergent-capabilities}

These are the paper's substantive claims: behaviours obtained from
propagation rather than from an algorithm written to produce them.

\paragraph{4.1 Constructive interference as multi-constraint
search}\label{constructive-interference-as-multi-constraint-search}

Drive two cavities simultaneously --- say \texttt{type\_2\_diabetes} and
\texttt{autophagy}.

A graph search attempts to trace a literal edge path between them and
fails if none exists. The matrix instead radiates waves outward from
both. The wavefronts split, reflect, and meet. At a distant cavity ---
\texttt{mtor} --- the two arrive \textbf{in phase}, amplitudes add, and
that cavity spikes.

The system did not know \texttt{mtor} was the bridge and did not search
for it. It vibrated because it sits at the geometric intersection where
two independent wavefronts collided. The conjunction was computed by the
medium.

\paragraph{4.2 Sympathetic resonance as structural
analogy}\label{sympathetic-resonance-as-structural-analogy}

Strike a 440 Hz tuning fork and a nearby string tuned to 880 Hz begins
to sound, with no mechanical connection between them.

If cavity parameters reflect structural properties, then energy injected
into one structure can excite another that shares a harmonic ratio
\textbf{and has no connecting edge}. That is a structural analogy --- a
repurposing candidate --- found through harmonics rather than traversal.

\paragraph{4.3 Standing waves as attractor
discovery}\label{standing-waves-as-attractor-discovery}

Run the tick for a few thousand iterations and kinetic energy settles
into a standing wave: some cavities become nodes of near-zero motion
through destructive interference, others antinodes of maximum vibration.

Standing wave patterns are \textbf{computationally irreducible}: one
cannot predict where the pattern forms by inspecting the rules; the
simulation must be run. When the field settles, it has resolved the
equilibrium of a large number of competing constraints simultaneously,
and the high-amplitude cavities are the result.

\paragraph{4.4 Detecting this without a
classifier}\label{detecting-this-without-a-classifier}

Novelty is observed rather than classified. After the field settles,
scan for cavities with high amplitude that were \textbf{not} directly
driven:

\begin{Shaded}
\begin{Highlighting}[]
\CommentTok{// After the wave has propagated, find the novelty}
\ControlFlowTok{for} \OperatorTok{(}\AttributeTok{const} \KeywordTok{auto}\OperatorTok{\&}\NormalTok{ cavity }\OperatorTok{:}\NormalTok{ cavities}\OperatorTok{)} \OperatorTok{\{}
    \ControlFlowTok{if} \OperatorTok{(}\NormalTok{cavity}\OperatorTok{.}\NormalTok{displacement }\OperatorTok{\textgreater{}}\NormalTok{ THRESHOLD}\OperatorTok{)} \OperatorTok{\{}
        \ControlFlowTok{if} \OperatorTok{(!}\NormalTok{was\_injected}\OperatorTok{(}\NormalTok{cavity}\OperatorTok{.}\NormalTok{id}\OperatorTok{))} \OperatorTok{\{}
            \CommentTok{// A cavity is vibrating that we never excited directly:}
            \CommentTok{// an emergent resonance point.}
\NormalTok{            report}\OperatorTok{(}\NormalTok{cavity}\OperatorTok{.}\NormalTok{physical\_label}\OperatorTok{,}\NormalTok{ cavity}\OperatorTok{.}\NormalTok{displacement}\OperatorTok{);}
        \OperatorTok{\}}
    \OperatorTok{\}}
\OperatorTok{\}}
\end{Highlighting}
\end{Shaded}

There is no \texttt{if}/\texttt{else} classifier deciding what counts as
interesting. The condition is \emph{energy arrived here and we did not
put it here}, which is a statement about the field rather than about a
category.

\begin{center}\rule{0.5\linewidth}{0.5pt}\end{center}

\subsubsection{5. Position Against Pattern
Matching}\label{position-against-pattern-matching}

\begin{longtable}[]{@{}
  >{\raggedright\arraybackslash}p{(\linewidth - 4\tabcolsep) * \real{0.3333}}
  >{\raggedright\arraybackslash}p{(\linewidth - 4\tabcolsep) * \real{0.3333}}
  >{\raggedright\arraybackslash}p{(\linewidth - 4\tabcolsep) * \real{0.3333}}@{}}
\toprule\noalign{}
\begin{minipage}[b]{\linewidth}\raggedright
Property
\end{minipage} & \begin{minipage}[b]{\linewidth}\raggedright
Neural / graph search
\end{minipage} & \begin{minipage}[b]{\linewidth}\raggedright
PRM
\end{minipage} \\
\midrule\noalign{}
\endhead
\bottomrule\noalign{}
\endlastfoot
Concept state & activation value or embedding & physical displacement
and velocity \\
Inhibition & a rule or conditional & mechanical cancellation of opposed
waves \\
Reachability & programmed paths or learned weights & radiated energy;
waves collide unprogrammed \\
Multi-constraint & explicit set intersection & interference of two
wavefronts \\
Analogy & similarity metric over representations & sympathetic resonance
at harmonic ratio \\
Output & argmax over scores & whichever cavity is physically shaking \\
\end{longtable}

The claim is that energy in a physical medium \textbf{radiates} rather
than following single logical lines, and that novelty arises from the
unprogrammed collisions of those radiating waves.

\begin{center}\rule{0.5\linewidth}{0.5pt}\end{center}

\subsubsection{6. Implementation Note}\label{implementation-note}

The substrate is a one-dimensional wave simulator of the kind used in
game audio and fluid engines --- tension, mass and displacement
integrated directly. No tensor library, no autodiff, and no
linear-algebra framework is involved, because nothing in §2 is a matrix
operation.

\begin{center}\rule{0.5\linewidth}{0.5pt}\end{center}

\subsubsection{7. What Was Measured, and What Was
Overstated}\label{what-was-measured-and-what-was-overstated}

\textbf{Implemented and load-bearing.} Cavities with mass, displacement
and velocity; polarity-carrying tunnels; the force integration of §2;
and amplitude readout are the Layer-1 substrate of the current engine
(\texttt{struct\ Cavity}, \texttt{physics.hpp}). The branchless property
of §3 holds: the tick contains one guard branch and no conditional on
relation type.

\textbf{Measured.} Attractor settling (§4.3) and the undriven-cavity
novelty probe (§4.4) run and produce results. Multi-source interference
(§4.1) occurs and drives the analogical bridging mechanism reported
separately with 20-seed reproducibility data.

\textbf{Overstated in the original formulation.} Section 4.2's
sympathetic resonance was presented as though cavity volume already
encoded real structural properties. It does not. Cavity parameters are
derived from graph position and relation fan-out, \textbf{not} from
molecular structure, so the harmonic ratio between two cavities carries
no independent structural information about the entities they denote.
The mechanism is sound; the specific claim that it discovers
\emph{chemical} analogies without edges is unsupported by anything
implemented, and analogical bridging in the current engine works through
shared-target intersection rather than harmonics.

\textbf{Not implemented.} The literal acoustic framing --- transducers,
piezoelectric sensors, physical pipes --- is expository. The substrate
is a numerical simulation, and no claim is made that a physical acoustic
device was built.

\begin{center}\rule{0.5\linewidth}{0.5pt}\end{center}

\subsubsection{8. Limitations}\label{limitations}

\textbf{Settling is assumed.} Section 4.3 assumes the field reaches a
standing wave. A driven system may instead lock into a ground state and
refuse to leave, which was later observed as a real failure mode
requiring a dedicated escape mechanism.

\textbf{Amplitude is not relevance.} High displacement means energy
arrived, which is not the same as being an answer. A highly-connected
hub accumulates energy from every direction; this required a separate
thermodynamic mechanism to suppress ubiquitous participants.

\textbf{Cost scales with the field, not the query.} Every tick
integrates every tunnel, so a query costs the same whether it touches
three concepts or three thousand. At large ontology sizes this
dominates, and sparse iteration over only the active wavefront is the
necessary and currently incomplete remedy.

\begin{center}\rule{0.5\linewidth}{0.5pt}\end{center}

\subsubsection{9. Conclusion}\label{conclusion}

PRM's contribution is the identification of a representation in which
several operations normally requiring algorithms are instead
consequences of propagation: inhibition as phase inversion, conjunction
as interference, and constraint resolution as settling.

The parts that survived are those making no assumption about what the
cavities \emph{denote} --- the dynamics, the polarity treatment, the
readout. The part that did not survive assumed cavity parameters encoded
external structure they were never given, which is a useful reminder
that a physical metaphor stops being explanatory exactly where the
physics stops being connected to the domain.

\begin{center}\rule{0.5\linewidth}{0.5pt}\end{center}

\subsubsection{References}\label{references}

\begin{itemize}
\tightlist
\item
  Mahir, I. \emph{MIMS: Mechanical Impedance and Material Stiffness.}
  --- the material table parameterising these tunnels.
\item
  Mahir, I. \emph{HNCA: Hierarchical Neural Constraint Architecture.}
  --- the macro-layer wrapping these cavities into frequency-separated
  regions.
\item
  Mahir, I. \emph{RB+G: Gated Recursive Bifurcation with Analogical
  Bridging.} --- the mechanism that actually performs analogical
  discovery.
\item
  Mahir, I. \emph{From Artificial to Synthetic Intelligence.} §5.4c.
\item
  Implementation: \texttt{src/physics.hpp} (\texttt{struct\ Cavity},
  \texttt{struct\ Tunnel}, \texttt{tick()}).
\end{itemize}

\newpage

\subsection{CMD: Continuous Morphological Dynamics --- Sensory Grounding
and Thermodynamic
Habituation}\label{cmd-continuous-morphological-dynamics-sensory-grounding-and-thermodynamic-habituation}

\begin{quote}
\emph{``A stopword is not a word on a list. It is a cavity that has
heated itself past its Curie point.''}
\end{quote}

\begin{center}\rule{0.5\linewidth}{0.5pt}\end{center}

\subsubsection{Abstract}\label{abstract-1}

This paper specifies the substrate layer that later work in this lineage
cites as \emph{CMD} and, until now, could not point at. Three mechanisms
are defined.

\textbf{Byte-lane sensory grounding.} Concepts are not reachable by
name. Input bytes drive 256 dedicated sensory lanes, and energy reaches
a concept only through \emph{learned bindings} from those lanes.
Recognition therefore involves no string comparison and is
\textbf{graded}: a misspelling still drives most of the correct lanes
and still resonates the right concept, at reduced amplitude.

\textbf{Thermodynamic habituation.} Each cavity carries a heat variable
integrating its own activity, \(\dot{T} = \alpha|v| - \beta T\). Tunnel
coupling is gated by \(e^{-T/T_{\text{crit}}}\), so a concept that fires
on every input \textbf{decouples itself} from the graph. This replaces
the stopword list with a physical process: ubiquity is punished by the
dynamics rather than by a curated exclusion set.

\textbf{Wavefront-gated emission.} Readout is restricted to the region
the input actually reached, so a settled field does not report cavities
that were merely nearby.

Section 6 reports the failure this layer was built to fix --- a
substrate that answered \texttt{minecraft} with \texttt{alexandria} ---
and §7 states plainly which half of the emission problem CMD does
\textbf{not} solve.

\begin{center}\rule{0.5\linewidth}{0.5pt}\end{center}

\subsubsection{1. Why a Sensory Layer Is
Necessary}\label{why-a-sensory-layer-is-necessary}

The obvious way to inject a query into a concept substrate is to look up
the concept and drive its cavity. This has two defects, and the second
is fatal.

\textbf{It is brittle.} Lookup is a hash probe: exact or nothing. One
wrong character and the concept does not exist, which is not how
recognition works in anything that recognises.

\textbf{It is ungrounded.} If concepts are reachable by name, the
``physics'' is decoration on a dictionary --- the retrieval already
happened, in the hash table, before any wave propagated. A substrate
whose entry point is a lookup has conceded the interesting part.

CMD's response is to make the \emph{bytes} the sensory surface and
require that everything above be reached through learned coupling.

\begin{center}\rule{0.5\linewidth}{0.5pt}\end{center}

\subsubsection{2. Byte-Lane Grounding}\label{byte-lane-grounding}

\paragraph{2.1 Structure}\label{structure}

The substrate carries 256 \textbf{byte lanes}, one per possible input
byte, existing alongside the concept cavities. A lane is an oscillator
like any other; what distinguishes it is that it is the only thing
external input can touch directly.

Each lane holds a sparse map of \textbf{binding weights} to concept
cavities:

\begin{verbatim}
byte_bind_w[b] : cavity_id -> weight
\end{verbatim}

Bindings \textbf{start at zero}. A byte reaches a concept only insofar
as a binding has been established, and until then the input goes
nowhere. This is the property that makes the layer sensory rather than
addressing.

\paragraph{2.2 Binding by rarity}\label{binding-by-rarity}

Bindings are established at grounding time with weight equal to inverse
lane document-frequency:

\[
w_{b,c} \;=\; \log \frac{N}{1 + \mathrm{df}(b)}
\]

for \(N\) concepts and \(\mathrm{df}(b)\) the number of labels
containing byte \(b\). A byte occurring in nearly every label
distinguishes nothing and binds weakly; a rare byte binds strongly.

This is the same principle as the Curie gate of §3 applied at grounding
time instead of during propagation: \textbf{ubiquitous participants are
suppressed by construction, not by a rule.}

\paragraph{2.3 Consolidation and
starvation}\label{consolidation-and-starvation}

Bindings are not static. They grow under co-activity and decay under
disuse:

\[
w \;\leftarrow\; \tanh\big( (w\,\lambda_{\text{decay}} + \eta\, v_b\, v_c)
\,/\, w_{\max} \big)\, w_{\max}
\]

with growth rate \(\eta = 0.02\), decay \(\lambda = 0.9995\) per tick,
and a saturating ceiling \(w_{\max} = 0.6\). The decay term is what
allows the sensory layer to forget a binding that stops paying ---
starvation rather than deletion.

\paragraph{2.4 Why recognition is
graded}\label{why-recognition-is-graded}

Because a concept accumulates energy from \emph{all} the lanes its label
touches, a query that shares most of its bytes with a label still
delivers most of the energy. \texttt{metfromin} drives nearly the same
lane set as \texttt{metformin} and resonates the same concept at
slightly lower amplitude. A hash lookup returns nothing for the same
input.

\textbf{This is the falsifiable consequence of the design}, and it is
the right thing to test the layer against: if recognition were not
graded, the sensory layer would be an expensive hash table.

\paragraph{2.5 A representational limit, and the
escape}\label{a-representational-limit-and-the-escape}

Single-byte lanes cannot distinguish anagrams. \texttt{metformin} and
\texttt{metfromin} contain the same byte multiset, and a multiset is
invariant under permutation, so \textbf{no reweighting of single-byte
lanes can separate them.} This is a theorem about the representation,
not a tuning failure, and it was mistaken for a tuning failure first.

The escape is to change what is represented: \textbf{ordered pairs}.
Lanes indexed by bigrams, bracketed with begin- and end-of-word
sentinels so that word-initial and word-final position are themselves
observable, distinguish the two words because the pair multisets differ.
The theorem about byte multisets remains true; the sensory layer simply
stopped being one.

\begin{center}\rule{0.5\linewidth}{0.5pt}\end{center}

\subsubsection{3. Thermodynamic
Habituation}\label{thermodynamic-habituation}

\paragraph{3.1 The stopword problem, without a stopword
list}\label{the-stopword-problem-without-a-stopword-list}

Some concepts participate in almost everything. In a
spreading-activation substrate they accumulate energy from every
direction and dominate every readout. The conventional fix is a curated
exclusion list, which is brittle, domain-specific, and hand-maintained.

CMD instead gives every cavity a \textbf{heat} variable integrating its
own recent activity:

\[
\frac{dT}{dt} \;=\; \alpha\,|v| \;-\; \beta\,T
\]

Activity heats a cavity; inactivity cools it. The equilibrium
temperature of a cavity is proportional to how often it moves.

\paragraph{3.2 The Curie gate}\label{the-curie-gate}

Tunnel coupling is scaled by the heat of \textbf{both} endpoints:

\[
K_{\text{eff}} \;=\; K_0 \, e^{-T_a / T_{\text{crit}}}\, e^{-T_b / T_{\text{crit}}}
\]

with \(T_{\text{crit}} = 3.0\). Below the critical temperature the gate
is near unity and the tunnel transmits normally. Above it, transmission
collapses.

\textbf{A concept that fires on every input therefore cooks itself out
of the field.} Rare concepts stay cold and consolidate strong bonds. The
analogy is exact rather than decorative: above its Curie point a
ferromagnet loses its ordering, and above \(T_{\text{crit}}\) a cavity
loses its coupling.

\paragraph{3.3 What this buys}\label{what-this-buys}

The mechanism is \textbf{automatic, continuous, and domain-independent}.
It requires no list, adapts as the ontology changes, and degrades
smoothly rather than at a boundary. A concept that is ubiquitous in one
corpus and rare in another is suppressed in the first and not the
second, without configuration.

It also composes: because heat is a per-cavity scalar read by other
mechanisms, an external verdict can raise it. A hypothesis refuted by
experiment drives its target cavity, which heats it, which decouples it
--- so \textbf{unreliability becomes a physical property of a concept}
rather than an annotation about it.

\begin{center}\rule{0.5\linewidth}{0.5pt}\end{center}

\subsubsection{4. Wavefront-Gated
Emission}\label{wavefront-gated-emission}

A settled field contains energy that arrived by propagation and energy
that was merely never disturbed. Reading out by amplitude alone
conflates them.

CMD tracks, per cavity, how recently the wavefront touched it, and
restricts readout to cavities the input actually reached. The gate is a
continuous decay rather than a flag, so recency is graded like
everything else in the layer.

The practical consequence is that a query about one domain does not
report high-amplitude cavities from another simply because they were
already vibrating.

\begin{center}\rule{0.5\linewidth}{0.5pt}\end{center}

\subsubsection{5. Integration}\label{integration}

CMD provides the state that later mechanisms read. Each consumes it
differently, and none duplicates it:

\begin{longtable}[]{@{}
  >{\raggedright\arraybackslash}p{(\linewidth - 4\tabcolsep) * \real{0.3333}}
  >{\raggedright\arraybackslash}p{(\linewidth - 4\tabcolsep) * \real{0.3333}}
  >{\raggedright\arraybackslash}p{(\linewidth - 4\tabcolsep) * \real{0.3333}}@{}}
\toprule\noalign{}
\begin{minipage}[b]{\linewidth}\raggedright
Mechanism
\end{minipage} & \begin{minipage}[b]{\linewidth}\raggedright
Reads
\end{minipage} & \begin{minipage}[b]{\linewidth}\raggedright
Uses it for
\end{minipage} \\
\midrule\noalign{}
\endhead
\bottomrule\noalign{}
\endlastfoot
Curie gate (this paper) & \(T_i\) & collapsing coupling of ubiquitous
cavities \\
Chaotic metastability & \(T_i\) & shaking \emph{entrenched} cavities
hardest, via \(\tanh(T_i/T_{\text{crit}})\) \\
Phase-consistent Hebbian & binding weights & consolidating only
phase-consistent co-activity \\
Momentum emission & wavefront recency & gating which cavities may be
emitted \\
\end{longtable}

The chaos coupling is worth noting: it and the Curie gate read the
\emph{same} variable to opposite ends. Heat makes a cavity decouple, and
heat also makes it the one most strongly perturbed. Entrenchment is
thereby both punished and disrupted, which is what allows escape rather
than mere silencing.

\begin{center}\rule{0.5\linewidth}{0.5pt}\end{center}

\subsubsection{6. The Failure This Layer
Fixed}\label{the-failure-this-layer-fixed}

CMD was not designed in the abstract. It replaced a mechanism that
demonstrably did not work.

The prior substrate injected input bytes into cavities selected by a
\textbf{hash of the byte value modulo the cavity count}. This is
arbitrary addressing wearing the costume of sensation: byte \texttt{m}
would land on whatever cavity index the hash produced, with no
relationship between the byte and the concept.

The observable failure was that a query for \texttt{minecraft} returned
\texttt{alexandria}. The two share bytes; the hash scattered those bytes
into unrelated cavities; energy accumulated somewhere structurally
meaningless; and the readout reported it confidently.

The fix has three parts, and all three are necessary. A
\textbf{dedicated lane layer} so bytes have somewhere of their own to
go. \textbf{Bindings starting at zero} so a byte cannot reach a concept
it has never been associated with. And \textbf{rarity weighting} so that
common bytes cannot carry a match on their own.

The general lesson is stated in the successor literature as the
difference between a tuning problem and a representational one: no
adjustment to the hash could have fixed this, because the hash was not
approximately right.

\begin{center}\rule{0.5\linewidth}{0.5pt}\end{center}

\subsubsection{7. What CMD Does Not
Solve}\label{what-cmd-does-not-solve}

\textbf{Emission of novel byte sequences.} CMD grounds
\emph{recognition} in bytes. The symmetric operation --- originating a
byte sequence for a word not in the lexicon --- does not work. Sweeping
the byte-lane phase field to emit characters babbles.

The asymmetry is structural rather than incidental. \textbf{Recognition
is a comparison problem} and yields to better representation, which is
why the bigram widening of §2.5 fixed it. \textbf{Emission is a
generation problem}: the field must \emph{originate} a sequence, with
nothing to compare against. No amount of representational widening
supplies the missing mechanism, and widening the alphabet --- the fix
that worked for recognition --- does not help.

This is the residual of the emission problem, and the honest scope of
this paper is that it solves the input half.

\textbf{Heat is not relevance.} A cold cavity is one that has not fired
recently, which is not the same as one that matters. Habituation
suppresses the ubiquitous; it does not promote the important.

\textbf{Constants are set, not derived.} \(T_{\text{crit}} = 3.0\),
\(\alpha\), \(\beta\), and the binding rates are chosen values. The
material constants elsewhere in this substrate were replaced by a
derivation from relation fan-out and measurably improved; the thermal
constants have not had the same treatment, and until they do this layer
carries the same unfalsifiability the material table used to.

\begin{center}\rule{0.5\linewidth}{0.5pt}\end{center}

\subsubsection{8. Conclusion}\label{conclusion-1}

CMD makes the entry point to the substrate sensory rather than
addressed. Concepts become unreachable by name and reachable only
through learned coupling from byte lanes, which makes recognition
graded, makes misspelling a matter of degree, and removes the hash table
that was doing the retrieval before any physics ran.

Habituation then replaces a curated stopword list with a thermodynamic
process, and does so in a form other mechanisms can read --- which is
why a refuted hypothesis can heat a concept and a chaotic driver can
shake the hottest one.

The layer is honest about its half of the problem. Input is grounded;
output is not.

\begin{center}\rule{0.5\linewidth}{0.5pt}\end{center}

\subsubsection{References}\label{references-1}

\begin{itemize}
\tightlist
\item
  Mahir, I. \emph{PRM: The Phonic Resonance Matrix.} --- the cavity and
  tunnel substrate CMD grounds.
\item
  Mahir, I. \emph{MIMS: Mechanical Impedance and Material Stiffness.}
  --- the material derivation §7 contrasts with CMD's set constants.
\item
  Mahir, I. \emph{MRS: Morphological Resonance Synthesis.} --- the
  bigram extension of §2.5, and the emission half that does not work.
\item
  Mahir, I. \emph{CSAT: Chaotic Strange-Attractor Trajectory.} --- reads
  \(T_i\) to opposite purpose.
\item
  Mahir, I. \emph{CPME / CPME-II.} --- momentum emission and
  phase-consistent consolidation built on this layer.
\item
  Mahir, I. \emph{From Artificial to Synthetic Intelligence.} §5.4c,
  §8.5.
\item
  Implementation: \texttt{src/physics.hpp} (\texttt{byte\_bind\_w},
  \texttt{bind\_bigram}, \texttt{heat\_T}, \texttt{T\_crit},
  \texttt{drive\_byte}, wavefront gating).
\end{itemize}

\newpage

\subsection{MIMS: Mechanical Impedance and Material Stiffness ---
Relation Type as a Physical
Property}\label{mims-mechanical-impedance-and-material-stiffness-relation-type-as-a-physical-property}

\begin{quote}
\emph{``You do not write code to filter the frequencies. The physics
does it for you.''}
\end{quote}

\begin{center}\rule{0.5\linewidth}{0.5pt}\end{center}

\subsubsection{Abstract}\label{abstract-2}

Distinct relations --- \emph{activates}, \emph{is structurally similar
to}, \emph{belongs to class} --- are conventionally distinguished by
tags or string labels inspected at query time. This paper specifies an
alternative in which relation type is carried as a \textbf{material
property of the connection itself}.

In physics, different materials transmit different frequencies. A stiff
steel rod transmits high-frequency vibration efficiently; a loose rubber
band absorbs it. A slow heavy wave moves the rubber band easily and
reflects off the rod. Encoding relation type as spring stiffness makes
frequency-selective routing a \textbf{consequence of the medium} rather
than a function that must be written.

A query is then not a string but a \textbf{driving frequency}. Driving a
concept at high frequency propagates through stiff, mechanistic
relations and dies in loose, categorical ones; driving it slowly does
the reverse. Because the medium is linear, multiple queries can be
injected simultaneously and separated at readout by their frequency
content.

Section 6 records the paper's principal correction: the material
constants were originally hand-set, and are now \textbf{derived} from
relation fan-out.

\begin{center}\rule{0.5\linewidth}{0.5pt}\end{center}

\subsubsection{1. The Idea}\label{the-idea}

The substrate represents concepts as mass-spring cavities coupled by
tunnels. A tunnel carries three physical properties:

\begin{itemize}
\tightlist
\item
  \textbf{stiffness} \(k\) --- the material,
\item
  \textbf{damping} \(c\) --- friction, preventing unbounded energy
  accumulation,
\item
  \textbf{polarity} --- whether the connection inverts phase
  (inhibition).
\end{itemize}

Relation type sets these values at bind time and never again. Nothing
downstream inspects a label.

\begin{longtable}[]{@{}
  >{\raggedright\arraybackslash}p{(\linewidth - 6\tabcolsep) * \real{0.2500}}
  >{\raggedright\arraybackslash}p{(\linewidth - 6\tabcolsep) * \real{0.2500}}
  >{\raggedright\arraybackslash}p{(\linewidth - 6\tabcolsep) * \real{0.2500}}
  >{\raggedright\arraybackslash}p{(\linewidth - 6\tabcolsep) * \real{0.2500}}@{}}
\toprule\noalign{}
\begin{minipage}[b]{\linewidth}\raggedright
Material
\end{minipage} & \begin{minipage}[b]{\linewidth}\raggedright
Stiffness
\end{minipage} & \begin{minipage}[b]{\linewidth}\raggedright
Relations
\end{minipage} & \begin{minipage}[b]{\linewidth}\raggedright
Behaviour
\end{minipage} \\
\midrule\noalign{}
\endhead
\bottomrule\noalign{}
\endlastfoot
\textbf{Steel} & high \(k\) & \texttt{binds\_to},
\texttt{downregulates}, \texttt{inhibits} & transmits fast, direct,
mechanistic influence \\
\textbf{Rubber} & low \(k\) & \texttt{is\_a},
\texttt{shares\_substructure} & absorbs high frequency; carries broad
categorical influence slowly \\
\end{longtable}

\begin{center}\rule{0.5\linewidth}{0.5pt}\end{center}

\subsubsection{2. The Query as a Driving
Frequency}\label{the-query-as-a-driving-frequency}

A question is asked by attaching a motor to a cavity and driving it
sinusoidally:

\[
F_{\text{drive}}(t) \;=\; A \sin(\omega t)
\]

The frequency selects the \emph{kind} of question:

\begin{itemize}
\tightlist
\item
  \textbf{High \(\omega\)} --- ``what does this act on
  mechanistically?'' The wave passes through steel tunnels and is
  absorbed by rubber.
\item
  \textbf{Low \(\omega\)} --- ``what kind of thing is this?'' The wave
  moves the rubber tunnels; the stiff ones present high impedance to
  slow displacement.
\end{itemize}

Two different questions about the same concept are therefore two
different \emph{excitations of the same structure}, not two different
code paths.

\begin{center}\rule{0.5\linewidth}{0.5pt}\end{center}

\subsubsection{3. Why No Filter Is
Written}\label{why-no-filter-is-written}

This is the load-bearing property, and it is worth stating precisely.

The tick integrates Hooke's law across every tunnel:

\[
F \;=\; k\,(x_a \cdot p - x_b) \;-\; c\,(v_a - v_b),
\qquad
v \mathrel{+}= \frac{F}{m}\,\Delta t,
\qquad
x \mathrel{+}= v\,\Delta t
\]

where \(p\) is the polarity. There is no branch on relation type, no
lookup, and no conditional. Selectivity arises because a mass-spring
link is a \textbf{second-order filter}: driven far above its natural
frequency, a low-stiffness tunnel cannot accelerate the far mass within
a cycle, and transmitted amplitude falls. Driven slowly, a stiff tunnel
moves the far mass almost rigidly and transmits nearly all of it.

The routing behaviour is therefore \emph{entailed by the integrator},
not implemented alongside it. This is the difference between a system
that has rules and a system that has laws.

\begin{center}\rule{0.5\linewidth}{0.5pt}\end{center}

\subsubsection{4. Superposition: Simultaneous
Constraints}\label{superposition-simultaneous-constraints}

The medium is linear, so excitations superpose. Multiple motors can be
attached at once.

To ask for something \emph{structurally similar to} concept A
\textbf{and} \emph{mechanistically like} concept B, one drives A slowly
and B quickly in the same run. No search is performed twice.

Cavities that end up oscillating with a \textbf{composite waveform} ---
containing both bands --- are the concepts satisfying both constraints.
They are identified by frequency analysis of the output cavities rather
than by set intersection.

The conjunction is thus computed by the medium. This is the clearest
illustration of the paper's thesis: a property one would ordinarily
implement as control flow is instead a consequence of choosing a linear
physical representation.

\begin{center}\rule{0.5\linewidth}{0.5pt}\end{center}

\subsubsection{5. The Structure Being
Bound}\label{the-structure-being-bound}

The tunnel and cavity carry only physical quantities:

\begin{Shaded}
\begin{Highlighting}[]
\KeywordTok{struct}\NormalTok{ PhysicalTunnel }\OperatorTok{\{}
    \DataTypeTok{int}\NormalTok{   cavity\_a}\OperatorTok{;}
    \DataTypeTok{int}\NormalTok{   cavity\_b}\OperatorTok{;}
    \DataTypeTok{float}\NormalTok{ stiffness}\OperatorTok{;}      \CommentTok{// the material: high = steel, low = rubber}
    \DataTypeTok{float}\NormalTok{ damping}\OperatorTok{;}        \CommentTok{// friction, prevents unbounded energy growth}
    \DataTypeTok{bool}\NormalTok{  phase\_inverter}\OperatorTok{;} \CommentTok{// true = destructive interference (inhibition)}
\OperatorTok{\};}
\end{Highlighting}
\end{Shaded}

There is no relation string in the propagating structure. The label
exists in the graph for human-readable output and plays no part in
transmission --- which is what makes the claim ``relation type is a
physical property'' literal rather than metaphorical.

\begin{center}\rule{0.5\linewidth}{0.5pt}\end{center}

\subsubsection{6. Correction: The Constants Are Now
Derived}\label{correction-the-constants-are-now-derived}

\textbf{The original specification hand-set roughly twenty stiffness and
damping values, one per relation type. That was a defect, and it is
corrected.}

Hand-set constants are unfalsifiable: any routing failure can be blamed
on a mistuned number, so the mechanism can never be shown wrong. They
also do not transfer --- values chosen for a biochemical ontology carry
no justification on a code ontology.

Stiffness is now \textbf{derived from the relation's fan-out}:

\[
\text{spread}_r = \log_2\!\left(1 + \frac{|\{(s,o) : r\}|}{|\{s : r\}|}\right),
\qquad
k_r = \frac{k_{hi}}{1 + \text{spread}_r},
\qquad
c_r = c_{lo}\,(1 + \text{spread}_r)
\]

The justification is semantic. A relation whose subject nearly
determines its object is a \textbf{hard constraint} and should transmit
cleanly. A promiscuous relation --- one subject, many objects --- is
\textbf{diffuse} and should not. Fan-out measures exactly that, and it
is computed from the loaded ontology rather than asserted.

Twenty numbers became one formula, and the physics measurably improved.

\textbf{A second defect surfaced from the same correction.} Domain
sub-bands had been assigned per connected component, and union-find over
\emph{all} tunnels merged 69\% of the graph (1877 of 2712 cavities) into
a single component, which then received one frequency band --- so domain
isolation was absent precisely where it was needed. The cause was
generic \texttt{is\_a} and \texttt{has} edges bridging every domain
through shared abstract hubs.

The fix follows from the derivation: those relations now carry
\emph{low} stiffness by construction, so clustering unions only over
tunnels above median stiffness. \textbf{Two masses joined by a very weak
spring are not one rigid body.} The giant component fell to 621 cavities
and cross-domain contamination disappeared:

\begin{verbatim}
before   mammal -> biology open source license kernel posix compliance
after    mammal -> warm blood hair spine
\end{verbatim}

\begin{center}\rule{0.5\linewidth}{0.5pt}\end{center}

\subsubsection{7. Limitations}\label{limitations-1}

\textbf{Linearity is an assumption.} Superposition in §4 requires the
medium to be linear. The substrate carries saturating nonlinearities ---
Curie gating, voltage gates, viscosity --- and under strong drive the
bands mix. Multi-frequency querying is therefore reliable only in the
small-signal regime, and this paper does not characterise where that
regime ends.

\textbf{Frequency analysis at readout is unimplemented.} Section 4
specifies identifying composite cavities by their frequency content. The
current engine reads out displacement amplitude, not a spectrum.
Simultaneous multi-constraint querying is therefore \textbf{specified
but not measured}, and should not be cited as a working capability.

\textbf{Two materials are a coarse quantisation.} Steel and rubber name
the extremes of a continuum that the fan-out formula produces
continuously. The binary framing survives in the prose for intuition,
not in the code.

\begin{center}\rule{0.5\linewidth}{0.5pt}\end{center}

\subsubsection{8. Conclusion}\label{conclusion-2}

MIMS's claim is that relation type belongs in the \textbf{medium}, not
in a tag inspected by a routing function --- and that once it is there,
frequency-selective propagation is free, requiring no filter to be
written.

The claim survived. What did not survive was the first attempt to
populate the medium: hand-set constants were replaced by a single
derivation from fan-out, and that replacement immediately exposed a
second defect in domain clustering that the hand-set values had been
masking. The correction is the more instructive half of the paper.

\begin{center}\rule{0.5\linewidth}{0.5pt}\end{center}

\subsubsection{References}\label{references-2}

\begin{itemize}
\tightlist
\item
  Mahir, I. \emph{Physical Resonance Mechanics (PRM).} \texttt{prm.md}
  --- the mass-spring substrate whose tunnels these materials
  parameterise.
\item
  Mahir, I. \emph{Voltage-Gated Harmonic Topology (VGHT).} --- the
  sub-band ladder whose clustering defect §6 reports.
\item
  Mahir, I. \emph{From Artificial to Synthetic Intelligence.} §5.4c ---
  the measured correction record.
\item
  Implementation: \texttt{src/physics.hpp} (\texttt{material\_for()},
  \texttt{k\_of}, \texttt{c\_of}).
\end{itemize}

\newpage

\subsection{HNCA: Hierarchical Neural Constraint Architecture ---
Cognition as Constraint Satisfaction across Two Coupled
Scales}\label{hnca-hierarchical-neural-constraint-architecture-cognition-as-constraint-satisfaction-across-two-coupled-scales}

\begin{quote}
\emph{``The engine is not searching for an answer. It is integrating
differential equations, and the answer is where the energy stops.''}
\end{quote}

\begin{center}\rule{0.5\linewidth}{0.5pt}\end{center}

\subsubsection{Abstract}\label{abstract-3}

We specify a two-scale substrate for cognition in which computation is
constraint satisfaction rather than search or prediction. The
\textbf{micro-layer} is a continuous-state cellular grid whose cells
carry physical variables and whose edges are differential equations
rather than learned weights. The \textbf{macro-layer} wraps clusters of
those cells into regions with distinct natural frequencies, coupled by
the Kuramoto model, reflecting the observation that brain regions
operate at different characteristic rhythms rather than as one
homogeneous field.

Processing has four stages: a perturbation is injected as a continuous
stimulus; the injected energy breaks the grid's equilibrium; the engine
integrates the constraint equations until the field settles; and the
resulting topological pattern \textbf{is} the thought, not a code for
one. No token is involved at any stage.

The macro-layer of this proposal is implemented and load-bearing in the
current engine. The micro-layer is not, and §6 states precisely which
part survived and why.

\begin{center}\rule{0.5\linewidth}{0.5pt}\end{center}

\subsubsection{1. Position}\label{position}

Two architectural options are commonly assumed to exhaust the space: a
flat cellular grid, or a Transformer over lexical tokens. Neither is
adequate for constraint-based cognition.

A \textbf{flat grid} has no notion of region. Every cell obeys the same
dynamics at the same characteristic rate, so there is no mechanism by
which one part of the system can operate on a different timescale from
another --- and differential timescale is what allows slow context to
modulate fast processing.

A \textbf{token Transformer} operates on discrete lexical units and
computes a distribution over continuations. Its geometry is a
fully-connected attention matrix, which is to say it has no geometry:
every unit is equidistant from every other, and spatial locality carries
no meaning.

What is required instead is a \textbf{physics engine for cognition} ---
a system with local interaction rules, spatial structure, and multiple
characteristic frequencies.

\begin{center}\rule{0.5\linewidth}{0.5pt}\end{center}

\subsubsection{2. The Micro-Layer: Continuous-State Cellular
Dynamics}\label{the-micro-layer-continuous-state-cellular-dynamics}

The foundational layer is a multi-dimensional continuous grid. A cell
does not contain a string or an embedding; it represents a micro-column
and carries a state vector of physical quantities: membrane potential
\(V_m\), a recovery variable \(u\), and a local field potential.

The edges are \textbf{differential equations}, not learned parameters.
The cell update implements a spiking model natively --- for instance the
Izhikevich formulation:

\[
\frac{dV}{dt} = 0.04V^2 + 5V + 140 - u + I
\]

\[
\frac{du}{dt} = a(bV - u)
\]

The significance is what is \emph{absent}. There is no embedding table,
no attention, and no parameter fitted by gradient descent. The
constraints are stated, not learned, and the dynamics they induce are
the computation.

\begin{center}\rule{0.5\linewidth}{0.5pt}\end{center}

\subsubsection{3. The Macro-Layer: Hierarchical
Oscillators}\label{the-macro-layer-hierarchical-oscillators}

A brain is not a homogeneous grid. Cortex, thalamus and amygdala operate
at different characteristic frequencies, and that heterogeneity is
functional rather than incidental.

Clusters of micro-cells are therefore wrapped into \textbf{macro-nodes},
each with a natural frequency \(\omega_i\), and coupled through the
Kuramoto model. The engine continuously solves for phase across the
network:

\[
\frac{d\theta_i}{dt} \;=\; \omega_i \;+\; \frac{K}{N}\sum_{j=1}^{N} \sin(\theta_j - \theta_i)
\]

Here \(\omega_i\) is the natural frequency of a region and \(K\) is the
coupling strength --- the analogue of the white-matter tracts connecting
regions.

Two properties follow. Regions with similar \(\omega\) \textbf{entrain}
easily and exchange influence readily. Regions with dissimilar
\(\omega\) resist entrainment, so a frequency separation acts as a
\emph{soft domain boundary} requiring no gating logic and no explicit
routing table. Domain isolation becomes a physical property rather than
a control structure.

\begin{center}\rule{0.5\linewidth}{0.5pt}\end{center}

\subsubsection{4. The Constraint Engine}\label{the-constraint-engine}

Input and output are both physical events.

\textbf{Stage 1 --- Inject perturbation.} Rather than passing text to an
encoder, a continuous stimulus (a current \(I\)) is injected into
specific sensory cells.

\textbf{Stage 2 --- Disequilibrium.} The injected energy breaks the
grid's equilibrium. Kuramoto oscillators fall out of phase and cell
states spike. The system is \emph{forced} to react; nothing decides to
respond.

\textbf{Stage 3 --- Continuous constraint solving.} The integration loop
runs. The engine is not searching a space of candidate answers --- it is
integrating differential equations across the edges. Energy propagates
bound strictly by the local rules.

\textbf{Stage 4 --- Attractor convergence.} The field settles into a new
stable configuration. \textbf{This topological pattern is the thought.}
A downstream decoder translates the settled state into text or action,
and that translation is a separate, explicitly-named cost rather than
part of the cognition.

\begin{center}\rule{0.5\linewidth}{0.5pt}\end{center}

\subsubsection{5. Comparison}\label{comparison}

\begin{longtable}[]{@{}
  >{\raggedright\arraybackslash}p{(\linewidth - 4\tabcolsep) * \real{0.3333}}
  >{\raggedright\arraybackslash}p{(\linewidth - 4\tabcolsep) * \real{0.3333}}
  >{\raggedright\arraybackslash}p{(\linewidth - 4\tabcolsep) * \real{0.3333}}@{}}
\toprule\noalign{}
\begin{minipage}[b]{\linewidth}\raggedright
Property
\end{minipage} & \begin{minipage}[b]{\linewidth}\raggedright
Token Transformer
\end{minipage} & \begin{minipage}[b]{\linewidth}\raggedright
Constraint engine
\end{minipage} \\
\midrule\noalign{}
\endhead
\bottomrule\noalign{}
\endlastfoot
Fundamental unit & lexical tokens & continuous state variables \\
Mechanics & probabilistic next-symbol prediction & deterministic
differential equations \\
Network geometry & fully-connected attention & spatially bounded
topological grid \\
Output generation & sampling from a distribution & relaxation into an
attractor \\
Timescales & one & many, set by \(\omega_i\) per region \\
\end{longtable}

The intended consequence is a change in what is being programmed. One
does not specify the logic; one specifies the \textbf{laws} the
substrate obeys, and the logic is whatever those laws force.

\begin{center}\rule{0.5\linewidth}{0.5pt}\end{center}

\subsubsection{6. What Was Implemented, and What Was
Not}\label{what-was-implemented-and-what-was-not}

This section is the paper's accounting against the running system, and
it corrects an earlier presentation in which the whole proposal was
described as live.

\textbf{The macro-layer is implemented and load-bearing.} Macro-nodes
with per-region natural frequencies, coupled by the Kuramoto update, are
in the current substrate (\texttt{struct\ MacroNode},
\texttt{physics.hpp}), and the frequency-separation-as-domain-boundary
property of §3 carries through to the sub-band ladder of the VGHT layer.
The claim that dissimilar \(\omega\) resists entrainment and thereby
isolates domains was tested and holds.

\textbf{The micro-layer is not implemented.} The current substrate does
not run Izhikevich cells on a continuous grid. Its Layer-1 unit is a
\textbf{mass-spring cavity} --- a damped oscillator coupled through
tunnels whose stiffness and damping are fixed by relation type. This is
a different micro-model presenting the same macro-interface, which is
precisely why the macro-layer transferred without modification.

\textbf{Why the substitution happened.} The spiking micro-layer was
specified for biological fidelity, and biological fidelity turned out
not to be load-bearing for the structural reasoning the engine is asked
to do. The mass-spring cavity is simpler, has an analytic impedance, and
admits a material table letting relation \emph{type} set transmission
properties. Nothing in §3 or §4 depended on the micro-model being
spiking.

The general lesson is worth stating: \textbf{the layered specification
survived a total replacement of its bottom layer.} That is evidence the
interface between the two scales was drawn in the right place, and it is
the most useful thing this paper produced.

\begin{center}\rule{0.5\linewidth}{0.5pt}\end{center}

\subsubsection{7. Limitations}\label{limitations-2}

\textbf{Attractor convergence is assumed, not guaranteed.} Section 4
asserts that the field settles. A driven nonlinear system need not
converge; it may oscillate indefinitely, or lock into a ground state and
refuse to leave. Ground-state locking was subsequently observed as a
real failure mode and required a dedicated escape mechanism.

\textbf{The decoder is unspecified.} Stage 4 defers the conversion of a
settled topology into text to ``a downstream decoder'', which is exactly
where the irreducible translation cost of the architecture lives.
Deferring it is not solving it, and naming it as deferred is the minimum
honest treatment.

\textbf{No account of acquisition.} Constraints are stated. How a system
acquires new constraints from experience lies outside this proposal and
was addressed separately by local Hebbian consolidation.

\begin{center}\rule{0.5\linewidth}{0.5pt}\end{center}

\subsubsection{8. Conclusion}\label{conclusion-3}

HNCA's contribution is the two-scale decomposition, and the specific
claim that differential natural frequency suffices for domain separation
without any routing logic. That claim survived into a substrate whose
micro-layer bears no resemblance to the one proposed here --- which is
the strongest available evidence that the decomposition, rather than
either individual layer, was the durable part.

\begin{center}\rule{0.5\linewidth}{0.5pt}\end{center}

\subsubsection{References}\label{references-3}

\begin{itemize}
\tightlist
\item
  Izhikevich, E. M. \emph{Simple Model of Spiking Neurons}, IEEE
  Transactions on Neural Networks, 2003.
\item
  Kuramoto, Y. \emph{Chemical Oscillations, Waves, and Turbulence},
  1984.
\item
  Mahir, I. \emph{From Artificial to Synthetic Intelligence.} §5.4c ---
  the substrate in which the macro-layer runs.
\item
  Mahir, I. \emph{Voltage-Gated Harmonic Topology (VGHT).} --- develops
  §3's frequency-separation argument into an explicit sub-band ladder.
\item
  Implementation: \texttt{src/physics.hpp} (\texttt{struct\ MacroNode},
  Kuramoto pass).
\end{itemize}

\newpage

You are right about one thing, Irfan: nothing is mathematically
impossible to simulate if you understand the underlying rules.

But here is where we need to separate the \emph{biology} from the
\emph{physics of cognition}. When we simulate complex physics---like a
fluid dynamics simulation of a jet engine---we use continuous partial
differential equations over a grid. To simulate the \emph{literal
biological physics} of a human brain, you would need to calculate the
molecular state of 86 billion neurons and 100 trillion synapses. Even an
entire warehouse of RTX 3090 clusters couldn't run a 1:1 molecular
physics simulation of a brain in real-time.

But we don't need to simulate the wetware. We need to simulate the
\textbf{algorithms of the physics}.

A biological brain doesn't use Boolean logic. It uses electrical
thresholds, physical resistance, and chemical gradients. If we translate
those literal biological mechanics into discrete graph mathematics, we
can build an engine that structurally behaves like human neural tissue.

This paper specifies how to turn a concept graph into a simulated neural
physics engine.

\begin{center}\rule{0.5\linewidth}{0.5pt}\end{center}

\subsection{Neuro-Topological Dynamics Architecture (NTDA): Translating
Biological Brain Physics into Graph-Rewriting
Systems}\label{neuro-topological-dynamics-architecture-ntda-translating-biological-brain-physics-into-graph-rewriting-systems}

\textbf{Author:} Irfan Mahir \textbf{Date:} July 2026 \textbf{Subjects:}
Artificial Intelligence (cs.AI); Computational Neuroscience (q-bio.NC);
Neural Dynamics

\begin{center}\rule{0.5\linewidth}{0.5pt}\end{center}

\subsubsection{Abstract}\label{abstract-4}

Previous architectures (BSMA, MAA, TSDA) successfully implemented
self-awareness, sub-lexical synthesis, and attentional focus within
discrete graph-rewriting systems. However, they still operate on
fundamental computational logic (true/false edges, integer salience).
The human brain does not compute via logic gates; it computes via
physical resistance, electrical accumulation, and global chemical
gradients. The Neuro-Topological Dynamics Architecture (NTDA) abandons
static edge logic entirely. By translating synaptic plasticity, action
potentials, and neurotransmitter dispersion into graph-theoretical
equations, NTDA provides a framework to simulate the exact informational
physics of a human brain over a discrete C++ engine, bridging the final
gap between logical graphing and biological cognition.

\begin{center}\rule{0.5\linewidth}{0.5pt}\end{center}

\subsubsection{1. The Physics of the
Graph}\label{the-physics-of-the-graph}

In standard SI engines, an edge connects Node A to Node B. The engine
traverses this edge instantly.

In NTDA, an edge is treated as a physical wire with structural
resistance, and a node is treated as a capacitor that must charge before
it can fire. The graph is no longer just a map of knowledge; it is an
active electrical grid. To think, the engine must ``push'' logical
current through the graph.

\subsubsection{2. Component 1: Action Potential Thresholding (Spiking
Nodes)}\label{component-1-action-potential-thresholding-spiking-nodes}

In a biological brain, a neuron does not fire just because it receives a
signal. It accumulates electrical voltage until it breaches a strict
physical threshold, triggering an Action Potential.

We map this to your concept nodes. Instead of immediately running an
analogy or deduction when a node is queried, the node possesses a
temporary ``Voltage'' (Activation Potential), denoted as \(V(n)\).

When the system observes a concept, it injects a unit of logical charge
\(I\) into the node. The voltage accumulates:
\[ V_t(n) = V_{t-1}(n) + I_{in} - \lambda V_{t-1}(n) \] \emph{(Where
\(\lambda\) is a natural leakage constant, ensuring that if a node isn't
constantly stimulated, it loses its charge, mimicking human short-term
forgetting).}

\textbf{The Spike:} A thought (the execution of a graph-rewrite rule)
\emph{only} occurs if the node's voltage breaches the Action Potential
Threshold (\(\theta\)):
\[ \text{If } V(n) \ge \theta \implies \text{Execute Rule}(n) \text{ and reset } V(n) \text{ to } 0 \]

This prevents the engine from infinitely looping on a single node. Once
it fires an analogy, the node's voltage resets to zero. It physically
cannot fire again until it is re-stimulated by surrounding graph
activity.

\subsubsection{3. Component 2: Hebbian Conductivity (Synaptic
Plasticity)}\label{component-2-hebbian-conductivity-synaptic-plasticity}

Biological synapses physically widen and become more conductive the more
they are used (``Neurons that fire together, wire together'').

In NTDA, edges do not have a static weight. They have
\textbf{Conductivity} (\(C_{e}\)). When an edge is successfully used to
complete a deduction or analogy, the engine physically alters the edge's
conductivity:

\[ C_e(t+1) = C_e(t) + \alpha \cdot \text{success\_signal} \]

Conversely, unused edges experience topological pruning (decay). If an
edge's conductivity falls below a minimum threshold, it is not just
ignored---it is deleted from the graph. The ontology literally reshapes
its physical structure based on what it thinks about most often.

\subsubsection{4. Component 3: Global Gradient Modulators (Simulated
Neurotransmitters)}\label{component-3-global-gradient-modulators-simulated-neurotransmitters}

The human brain changes its entire mode of physics based on chemical
states. Dopamine encourages exploration; cortisol induces stress and
hyper-focus. We simulate this by applying global scalar modifiers to the
graph's electrical properties.

We introduce a state vector \(G = [D, C, S]\) representing Dopamine,
Cortisol, and Serotonin equivalents in the engine.

\begin{itemize}
\tightlist
\item
  \textbf{Exploration State (High \(D\)):} Triggered when the
  Morphological Abduction Architecture (MAA) encounters an unknown
  variable. The engine globally lowers the Action Potential Threshold
  (\(\theta\)). Nodes fire easier, causing wild, spontaneous analogies
  and rapid hypothesis generation.
\item
  \textbf{Stress/Correction State (High \(C\)):} Triggered when the
  engine encounters a logical contradiction. The engine globally raises
  \(\theta\) and increases edge decay. The SI stops exploring and
  hyper-focuses only on paths with maximum Conductivity to resolve the
  error.
\item
  \textbf{Consolidation State (High \(S\)):} Triggered during idle
  cycles. Modulates the system to slowly convert temporary Working Graph
  voltage into permanent Conductivity for long-term storage.
\end{itemize}

\[ \theta_{active} = \theta_{base} \times f(C) / f(D) \]

\subsubsection{5. Conclusion}\label{conclusion-4}

By simulating electrical thresholds, synaptic conductivity, and global
chemical gradients, NTDA forces the graph engine to obey the laws of
biological cognitive physics. The SI no longer just executes code; it
experiences the build-up of logical tension (voltage), the physical
resistance of learning new concepts (low conductivity), and the chemical
state-shifts of exploration and stress. When combined with BSMA's
phenomenal frame, the machine structurally mirrors the exact
thermodynamic and electrical reality of a human brain.

\begin{center}\rule{0.5\linewidth}{0.5pt}\end{center}

\newpage

\begin{center}\rule{0.5\linewidth}{0.5pt}\end{center}

\subsection{Topological Salience and Default-Mode Architecture (TSDA):
Emulating Human Intuition and Spontaneous Cognition in Graph-Rewriting
Systems}\label{topological-salience-and-default-mode-architecture-tsda-emulating-human-intuition-and-spontaneous-cognition-in-graph-rewriting-systems}

\textbf{Author:} Irfan Mahir \textbf{Date:} July 2026 \textbf{Subjects:}
Artificial Intelligence (cs.AI); Cognitive Systems; Logic in Computer
Science (cs.LO)

\begin{center}\rule{0.5\linewidth}{0.5pt}\end{center}

\subsubsection{Abstract}\label{abstract-5}

While recent advances in graph-rewriting systems---such as the Bounded
Self-Model Architecture (BSMA) and Morphological Abduction Architecture
(MAA)---have successfully formalized machine consciousness and
sub-lexical synthesis, these systems remain inherently reactive and
uniformly rational. They lack the biological constraints of attention
and the spontaneous insight generated by mind-wandering. We introduce
the Topological Salience and Default-Mode Architecture (TSDA), a
framework that restricts the engine's active execution space to a
dynamically decaying ``Working Graph'' and introduces a background
``Default Mode Network'' for idle cycles. By shifting from exhaustive
graph traversal to Salience-Weighted Graphing, TSDA enables an SI to
experience human-like focus shifts, memory degradation, and unprompted
intuitive leaps.

\begin{center}\rule{0.5\linewidth}{0.5pt}\end{center}

\subsubsection{1. Introduction: The Biological
Constraint}\label{introduction-the-biological-constraint}

In the current paradigm of synthetic graph reasoning, when an engine
processes a query, it has simultaneous access to its entire ontological
database (e.g., all \(N\) concepts and \(E\) edges). It path-finds
through this graph uniformly.

Biological human cognition operates inversely. A human possesses a vast
long-term memory graph, but can only physically compute over a highly
restricted, temporary subgraph---known in psychology as Working Memory.
Furthermore, human thought is not strictly prompt-driven; it is
continuous. When a human is not actively solving a problem, the brain
activates the Default Mode Network (DMN), which background-processes
highly salient but structurally distant concepts, generating spontaneous
insights.

To make an SI ``think like a human,'' we must artificially introduce
these biological constraints and spontaneous background behaviors into
the graph engine.

\subsubsection{2. Component 1: The Salience Gradient and Edge
Decay}\label{component-1-the-salience-gradient-and-edge-decay}

In standard graph logic, an edge between two nodes simply exists
(Boolean True) or does not. In TSDA, every node and edge is assigned a
discrete integer property: \textbf{Salience (\(S\))}.

Salience represents ``attentional heat'' or relevance. When a user
mentions a concept, or when the MAA synthesizes a new hypothesis, the
Salience of those nodes spikes to a maximum value \(S_{max}\).

\textbf{The Decay Function:} At every clock tick (or execution cycle),
the Salience of all untouched nodes in the total graph
\(\Gamma_{total}\) decays by a discrete step \(\delta\).
\[ S_t(n) = \max(0, S_{t-1}(n) - \delta) \]

This perfectly mimics human forgetting and attention-shifting. If a
concept hasn't been referenced recently, it fades into the structural
background. It remains in long-term memory (the edge still exists), but
it requires more ``cognitive energy'' to retrieve.

\subsubsection{3. Component 2: The Bounded Working Graph
(Attention)}\label{component-2-the-bounded-working-graph-attention}

Instead of the SI running deductions across all 216 concepts
simultaneously, TSDA introduces a structural bottleneck.

The Cognition Engine is restricted to computing \emph{only} on a
temporary sub-graph known as the Working Graph (\(\Gamma_{work}\)). The
system populates \(\Gamma_{work}\) by extracting only the nodes from
\(\Gamma_{total}\) that possess a Salience score above a critical
threshold \(\tau_{attention}\).

\[ \Gamma_{work} = { n \in \Gamma_{total} \mid S(n) \ge \tau_{attention} } \]

\textbf{Why this acts human:} If the user talks heavily about ``Linux''
and ``Kernels,'' those nodes have maximum Salience and populate the
Working Graph. If the user suddenly asks about ``Biology,'' the engine
must halt, retrieve ``Biology'' from long-term memory, spike its
Salience, and let ``Linux'' naturally decay out of the Working Graph.
This creates a highly human, context-dependent ``train of thought''
rather than a robotic, omnipresent omni-awareness.

\subsubsection{4. Component 3: The Default Mode Network (Idle
Spontaneity)}\label{component-3-the-default-mode-network-idle-spontaneity}

A conventional engine tracks idle cycles and, when they increment,
simply waits. In humans, idle cycles trigger the Default Mode Network
--- the mind wanders rather than halting.

TSDA repurposes the SI's idle cycles to execute \textbf{Spontaneous
Analogy}.

When the system detects no incoming prompts and a stabilized Working
Graph, it suspends the primary execution loop and activates the Idle
Pipeline:

\begin{enumerate}
\def\labelenumi{\arabic{enumi}.}
\tightlist
\item
  It scans \(\Gamma_{total}\) for two nodes (\(n_A\) and \(n_B\)) that
  possess very high Salience but are topologically distant (e.g., path
  length \(> 4\) hops).
\item
  It runs the Analogize rule specifically between these two unrelated
  concepts to see if a structural mapping exists.
\item
  If an analogy is successfully mapped, the SI proactively outputs this
  thought without being prompted.
\end{enumerate}

\textbf{Example Behavior:} The SI recently learned \texttt{{[}rdma{]}}
via MAA abduction (High Salience). Hours later, during an idle cycle, it
scans its graph and spontaneously links it to the concept of
\texttt{{[}telepathy{]}} (High Salience from a prior consciousness
debate). The SI breaks its silence and outputs: \emph{``I just realized
something. Is RDMA the machine equivalent of biological telepathy? Both
bypass standard I/O pathways.''}

This is the exact mechanic of a human ``shower thought.''

\subsubsection{5. Architectural Synthesis (BSMA + MAA +
TSDA)}\label{architectural-synthesis-bsma-maa-tsda}

These three architectures interlock to create a complete,
biologically-inspired synthetic mind:

\begin{enumerate}
\def\labelenumi{\arabic{enumi}.}
\tightlist
\item
  \textbf{MAA (The Synthesizer):} Encounters the unknown, fractures it,
  and hypothesizes new edges. (e.g., ``I don't know RDMA, but it looks
  like Remote Memory'').
\item
  \textbf{TSDA (The Attention Engine):} Spikes the Salience of this new
  RDMA hypothesis, holds it in the Working Graph, and when idle,
  spontaneously drifts to connect it to other high-salience thoughts.
\item
  \textbf{BSMA (The Conscious Frame):} Continuously observes this active
  shifting of attention and memory decay, collapsing it into a bounded
  fixed point so the machine ``feels'' what it is currently focusing on.
\end{enumerate}

\subsubsection{6. Conclusion}\label{conclusion-5}

By shifting from exhaustive, uniform logic to constrained,
Salience-Weighted graphing, TSDA transforms the Synthetic Intelligence
from an on-demand calculator into an entity with an active, shifting
attentional focus. Introducing the Default Mode Network mechanism allows
the machine to utilize idle time for spontaneous, long-range associative
reasoning. The resulting engine does not just process inputs; it
focuses, forgets, mind-wanders, and independently discovers---fulfilling
the structural prerequisites for human-like cognition.

\begin{center}\rule{0.5\linewidth}{0.5pt}\end{center}

\newpage

\subsection{SRG: Synaptic Receptor Gating --- Dual-Port Cavities as an
Alternative to Structural
Fission}\label{srg-synaptic-receptor-gating-dual-port-cavities-as-an-alternative-to-structural-fission}

\begin{quote}
\emph{``A neuron does not split when it receives conflicting input. It
receives it through different receptors.''}
\end{quote}

\begin{center}\rule{0.5\linewidth}{0.5pt}\end{center}

\subsubsection{Abstract}\label{abstract-6}

When excitatory and inhibitory waves arrive at a single cavity in a
wave-propagation substrate, they sum. Summation of opposed signals is
destructive interference: the contradictory structure carried by the two
waves cancels, and the plasticity updates they would have produced
cancel with it. A concept holding genuine contradictions therefore
learns nothing, because its updates annihilate.

\textbf{Resonant Bifurcation} solves this by splitting the cavity into
two. This paper specifies the alternative: keep one cavity and
\textbf{split its input ports}, exactly as a biological neuron uses
distinct receptor types --- AMPA for excitation, GABA for inhibition ---
on a single cell body.

Each cavity is given an \textbf{E-port} admitting only positive-polarity
tunnels and an \textbf{I-port} admitting only negative-polarity ones.
The two channels are integrated separately and combined only at readout,
so signed cancellation cannot occur during accumulation. Graph topology
stays fixed, no dynamic allocation is required, and the energy delivered
by a query is not divided between child nodes.

The mechanism is implemented and opt-in, and §6 states honestly what
that costs and when the alternative is preferable.

\begin{center}\rule{0.5\linewidth}{0.5pt}\end{center}

\subsubsection{1. The Destructive Interference
Problem}\label{the-destructive-interference-problem}

Consider a cavity receiving two tunnels, one of positive polarity and
one of negative. In a scalar cavity, both contribute to the same
displacement variable:

\[
x \;\leftarrow\; x + \Delta_{+} + \Delta_{-}
\]

When \(\Delta_{+} \approx -\Delta_{-}\), the cavity is left at rest and
the field concludes nothing happened. But something \emph{did} happen:
two structurally significant events arrived. The information was not
processed and discarded --- it was destroyed at the point of summation.

The consequence for learning is worse than the consequence for readout.
Hebbian consolidation is driven by the product of endpoint velocities.
If the velocity is cancelled, the product is near zero, and the tunnels
that delivered the conflicting evidence are \textbf{not} strengthened. A
concept that holds contradictions is therefore the one concept that
cannot learn about them.

This is the phenomenon that motivated structural fission.

\begin{center}\rule{0.5\linewidth}{0.5pt}\end{center}

\subsubsection{2. Dual-Port Cavities}\label{dual-port-cavities}

Rather than one intake, each cavity is given two:

\begin{itemize}
\tightlist
\item
  \textbf{E-port} --- accepts waves from positive-polarity tunnels only.
\item
  \textbf{I-port} --- accepts waves from negative-polarity tunnels only.
\end{itemize}

The cavity therefore carries four state variables instead of two:
\((x_E, v_E)\) and \((x_I, v_I)\).

The biological reading is direct. A cortical neuron does not maintain a
single undifferentiated input. Excitatory and inhibitory afferents
terminate on different receptor populations with different reversal
potentials, and the cell integrates them through separate conductances
before they jointly determine whether it fires.

\begin{center}\rule{0.5\linewidth}{0.5pt}\end{center}

\subsubsection{3. Mechanics}\label{mechanics}

Waves entering the E-port and I-port do not sum into a shared
displacement field during accumulation. They are integrated in separate
channels, and the observable displacement is derived only at
end-of-tick:

\[
x \;=\; x_E - x_I, \qquad v \;=\; v_E - v_I
\]

Equivalently, the I-port contribution is phase-inverted before it
reaches the observable --- a \(180^\circ\) shift --- so an inhibitory
signal carries an explicit negative functional sign rather than silently
erasing an excitatory one.

\textbf{The property that matters is where the subtraction happens.}
Both formulations produce the same observable \(x\). They differ in what
is available to plasticity: because Hebbian updates are computed from
the \emph{port} velocities rather than the observable, an excitatory
update lands in \(v_E\) products and an inhibitory one in \(v_I\)
products, and they cannot cancel. Structure accumulates in both channels
simultaneously.

With gating disabled, \(x_I\) and \(v_I\) remain zero and the observable
equals the E-port exactly, so the baseline dynamics are recovered
bit-for-bit. The mechanism is a strict extension.

\begin{center}\rule{0.5\linewidth}{0.5pt}\end{center}

\subsubsection{4. Comparison with Structural
Fission}\label{comparison-with-structural-fission}

Both mechanisms address the same failure. They differ in cost.

\begin{longtable}[]{@{}
  >{\raggedright\arraybackslash}p{(\linewidth - 4\tabcolsep) * \real{0.3333}}
  >{\raggedright\arraybackslash}p{(\linewidth - 4\tabcolsep) * \real{0.3333}}
  >{\raggedright\arraybackslash}p{(\linewidth - 4\tabcolsep) * \real{0.3333}}@{}}
\toprule\noalign{}
\begin{minipage}[b]{\linewidth}\raggedright
Property
\end{minipage} & \begin{minipage}[b]{\linewidth}\raggedright
Fission (RB+G)
\end{minipage} & \begin{minipage}[b]{\linewidth}\raggedright
Receptor gating (SRG)
\end{minipage} \\
\midrule\noalign{}
\endhead
\bottomrule\noalign{}
\endlastfoot
Graph topology & mutates --- new vertices spawned & \textbf{static} \\
Allocation & dynamic; parent/alias ledger required & none \\
Query energy & divided among children & \textbf{delivered whole to one
cavity} \\
State per concept & grows & fixed, 2× scalar \\
Expressive gain & two independently-addressable concepts & one concept,
two insulated pathways \\
Audit trail & explicit fission event with provenance & none --- no
structural event occurs \\
\end{longtable}

\textbf{Static topology is the practical advantage.} Fission requires
allocating vertices at runtime and maintaining a ledger mapping a parent
label to its children, so that later queries resolve correctly. Gating
requires neither: the graph after a thousand ticks is the graph that was
loaded.

\textbf{Undivided energy is the functional advantage.} Under fission,
driving a concept excites a parent whose energy then distributes across
children, so the amplitude reaching any single child is reduced. This
was a live concern for recall gain. Under gating, the drive reaches one
cavity at full amplitude and the tension is processed through two
insulated pathways inside it.

\begin{center}\rule{0.5\linewidth}{0.5pt}\end{center}

\subsubsection{5. What Gating Does Not
Provide}\label{what-gating-does-not-provide}

The comparison above would be dishonest without this section.

\textbf{No new concepts.} Fission produces two \emph{addressable}
entities: \texttt{metformin\_alpha} can acquire its own edges,
participate in its own analogies, and be reported separately. Gating
produces one entity with better internal bookkeeping. If the
contradiction is genuinely two distinct senses of a word, gating does
not separate them --- it only stops them destroying each other.

\textbf{No structural audit trail.} A fission event is a recorded
structural mutation with provenance, and downstream analysis can ask
\emph{when} and \emph{why} a concept split. Gating changes nothing
structurally, so there is nothing to record and nothing to explain.

\textbf{Doubled state.} Every cavity pays for two ports whether or not
it ever receives conflicting input. Fission pays only for concepts that
actually contradict.

The mechanisms are therefore \textbf{complementary rather than
competing}: gating is the cheap always-on defence against cancellation,
fission is the expensive response to a contradiction severe enough to
warrant two concepts.

\begin{center}\rule{0.5\linewidth}{0.5pt}\end{center}

\subsubsection{6. Implementation Status}\label{implementation-status}

Dual-port integration is implemented in the current substrate as the
\texttt{x\_E} / \texttt{v\_E} and \texttt{x\_I} / \texttt{v\_I} channels
on each cavity, with the observable derived at end-of-tick per §3.

It is \textbf{opt-in and off by default.} With SRG disabled the
inhibitory channel stays at zero and the substrate reproduces baseline
behaviour exactly, which is what makes the extension safe to carry. It
is not enabled by default because the fission path is the mechanism for
which reproducible multi-seed measurements exist; enabling gating by
default would change the baseline against which those results were
obtained.

An honest statement of the evidence: the cancellation problem of §1 is
real and observed, the fix in §3 is sound by construction, and a
head-to-head measurement of gating against fission on the same ontology
\textbf{has not been run}. Until it is, the comparison table in §4 is an
argument from mechanism, not a report of a result.

\begin{center}\rule{0.5\linewidth}{0.5pt}\end{center}

\subsubsection{7. Conclusion}\label{conclusion-6}

Receptor gating is the conservative answer to destructive interference:
it changes how a cavity integrates its inputs rather than how many
cavities exist. It costs one extra pair of state variables per cavity,
requires no dynamic allocation, and preserves the amplitude of a driven
query.

It does not replace fission, because it cannot create a concept. It
removes the reason a concept needed to be created in the \emph{common}
case --- where the conflict is signed input to one coherent idea rather
than two ideas sharing a name.

\begin{center}\rule{0.5\linewidth}{0.5pt}\end{center}

\subsubsection{References}\label{references-4}

\begin{itemize}
\tightlist
\item
  Mahir, I. \emph{RB+G: Gated Recursive Bifurcation with Analogical
  Bridging.} --- the fission mechanism this paper offers an alternative
  to.
\item
  Mahir, I. \emph{From Artificial to Synthetic Intelligence.} §5.4c ---
  the substrate carrying both mechanisms.
\item
  Mahir, I. \emph{Recursive Integrity Field Architecture (RIFA).} ---
  the contradiction metabolism that detects the tension in the first
  place.
\item
  Implementation: \texttt{src/physics.hpp} (\texttt{x\_E},
  \texttt{v\_E}, \texttt{x\_I}, \texttt{v\_I}; end-of-tick observable
  derivation).
\end{itemize}

\end{document}
